% landingpages-figure-system.tex
% Canonical shared TikZ visual system for technical landing-page figures.
%
% This file is the single cross-project style authority for newly migrated
% landing-page figures. It combines the reusable TikZ vocabulary originally
% developed for MQTTSuite with the authoritative SNode.C presentation contract.
%
% Presentation authority:
%   SNodeC/SNode.C-Book/assets/figures/src/snodec-figure-style.tex
%
% Canonical source rule:
% - figure .tex sources + this file are authoritative;
% - generated SVGs are publication output and are never hand-edited;
% - real runtime/UI evidence remains real raster capture.

\usetikzlibrary{
  arrows.meta,
  backgrounds,
  calc,
  fit,
  matrix,
  positioning,
  shapes.geometric,
  shapes.misc,
  shadows.blur
}

% ---------------------------------------------------------------------------
% Canonical palette
% ---------------------------------------------------------------------------

\definecolor{snodecInk}{HTML}{17212B}
\definecolor{snodecMuted}{HTML}{4B5F68}
\definecolor{snodecRule}{HTML}{7A8B93}
\definecolor{snodecBlue}{HTML}{0B4F6C}
\definecolor{snodecBlueSoft}{HTML}{E3F2F8}
\definecolor{snodecGreen}{HTML}{1F7A68}
\definecolor{snodecGreenSoft}{HTML}{E5F4EF}
\definecolor{snodecOrange}{HTML}{C46A1B}
\definecolor{snodecOrangeSoft}{HTML}{FAEFE3}
\definecolor{snodecRed}{HTML}{B24A4A}
\definecolor{snodecRedSoft}{HTML}{F8E8E8}
\definecolor{snodecGraySoft}{HTML}{F1F3F2}

% Compatibility names used by the shared semantic TikZ vocabulary.
\colorlet{MQText}{snodecInk}
\colorlet{MQTextMuted}{snodecMuted}
\colorlet{MQLine}{snodecMuted!76}
\colorlet{MQBorder}{snodecRule!85}
\colorlet{MQSurface}{white}
\colorlet{MQSurfaceAlt}{snodecGraySoft}
\colorlet{MQWhite}{white}
\colorlet{MQDecisionFill}{snodecOrangeSoft}
\colorlet{MQDecisionLine}{snodecOrange!80}
\colorlet{MQSuccessFill}{snodecGreenSoft}
\colorlet{MQSuccessLine}{snodecGreen!78}
\colorlet{MQErrorFill}{snodecRedSoft}
\colorlet{MQErrorLine}{snodecRed!82}
\colorlet{MQControlFill}{snodecBlueSoft}
\colorlet{MQControlLine}{snodecBlue!78}
\colorlet{MQWarningFill}{snodecOrangeSoft}
\colorlet{MQWarningLine}{snodecOrange!80}
\colorlet{MQCrossLane}{snodecBlue!70}
\colorlet{MQDataPlaneFill}{snodecBlueSoft}
\colorlet{MQDataPlaneLine}{snodecBlue!78}
\colorlet{MQAdminPlaneFill}{snodecGraySoft}
\colorlet{MQAdminPlaneLine}{snodecRule!85}
\colorlet{MQTrustFill}{white}
\colorlet{MQTrustLine}{snodecRule!82}
\colorlet{MQEvidenceRuntimeFill}{snodecGreenSoft}
\colorlet{MQEvidenceRuntimeLine}{snodecGreen!78}
\colorlet{MQEvidenceSourceFill}{snodecBlueSoft}
\colorlet{MQEvidenceSourceLine}{snodecBlue!78}

% ---------------------------------------------------------------------------
% Canonical optical systems
% ---------------------------------------------------------------------------
% These values are deliberately independent of the spatial layout rhythm. The
% root PAGE-SYSTEM.md keeps typography, strokes, arrowheads and radii in separate
% canonical optical systems.

\newcommand{\MQRadiusXS}{2.5pt}
\newcommand{\MQRadiusS}{2.5pt}
\newcommand{\MQRadiusM}{2.5pt}
\newcommand{\MQRadiusL}{3pt}
\newcommand{\MQRadiusXL}{3pt}

\newcommand{\MQLineWidthPrimary}{.70pt}
\newcommand{\MQLineWidthSecondary}{.58pt}
\newcommand{\MQArrowPrimaryLength}{2.2mm}
\newcommand{\MQArrowPrimaryWidth}{1.6mm}
\newcommand{\MQArrowSecondaryLength}{2.0mm}
\newcommand{\MQArrowSecondaryWidth}{1.4mm}

% ---------------------------------------------------------------------------
% Canonical spatial layout rhythm
% ---------------------------------------------------------------------------
% R is the one independently specified technical-figure layout distance.
% Every other spatial value below is calculated from R. Figure sources consume
% semantic tokens; they do not own millimetre values or local multipliers.

\newcommand{\MQLayoutRhythm}{8.3mm}

% Default rational rhythm scale from PAGE-SYSTEM.md.
\pgfmathsetlengthmacro{\MQRhythmMicro}{1/4*\MQLayoutRhythm}
\pgfmathsetlengthmacro{\MQRhythmCompact}{1/2*\MQLayoutRhythm}
\pgfmathsetlengthmacro{\MQRhythmTight}{3/4*\MQLayoutRhythm}
\pgfmathsetlengthmacro{\MQRhythmNormal}{\MQLayoutRhythm}
\pgfmathsetlengthmacro{\MQRhythmLoose}{4/3*\MQLayoutRhythm}
\pgfmathsetlengthmacro{\MQRhythmLarge}{5/3*\MQLayoutRhythm}

% Generic spacing compatibility names.
\newcommand{\MQU}{\MQRhythmMicro}
\newcommand{\MQGapXS}{\MQRhythmTight}
\newcommand{\MQGapS}{\MQRhythmTight}
\newcommand{\MQGapM}{\MQRhythmNormal}
\newcommand{\MQGapL}{\MQRhythmLoose}
\newcommand{\MQGapXL}{\MQRhythmLarge}

% Node interior and small label/evidence padding.
\pgfmathsetlengthmacro{\MQNodePadX}{3/8*\MQLayoutRhythm}
\pgfmathsetlengthmacro{\MQNodePadY}{3/10*\MQLayoutRhythm}
\pgfmathsetlengthmacro{\MQDecisionInnerPad}{1/20*\MQLayoutRhythm}
\pgfmathsetlengthmacro{\MQLabelInnerPad}{1/24*\MQLayoutRhythm}
\pgfmathsetlengthmacro{\MQEvidencePadX}{1/6*\MQLayoutRhythm}
\pgfmathsetlengthmacro{\MQEvidencePadY}{1/12*\MQLayoutRhythm}

% Canonical horizontal rhythm -- edge-to-edge gaps.
\newcommand{\MQHorizontalGapTight}{\MQRhythmTight}
\newcommand{\MQHorizontalGap}{\MQRhythmNormal}
\newcommand{\MQHorizontalGapLoose}{\MQRhythmLoose}
\newcommand{\MQHorizontalMilestoneGap}{\MQRhythmLarge}

% Vertical / branch rhythm and structural padding.
\newcommand{\MQBranchGapTight}{\MQRhythmTight}
\newcommand{\MQBranchGap}{\MQRhythmNormal}
\newcommand{\MQLaneGap}{\MQRhythmLarge}
\newcommand{\MQLanePadX}{\MQRhythmCompact}
\newcommand{\MQLanePadY}{\MQRhythmCompact}
\newcommand{\MQGroupPad}{\MQRhythmCompact}
\newcommand{\MQTrustPad}{\MQRhythmCompact}
\newcommand{\MQTitleGap}{\MQRhythmCompact}
\newcommand{\MQLabelGap}{\MQRhythmMicro}
\pgfmathsetlengthmacro{\MQFigureSafeX}{5/4*\MQLayoutRhythm}
\pgfmathsetlengthmacro{\MQFigureSafeY}{5/4*\MQLayoutRhythm}
\newcommand{\MQDesktopSafe}{\MQRhythmNormal}
\newcommand{\MQMobileSafe}{\MQRhythmNormal}
\newcommand{\MQMobileHorizontalGap}{\MQRhythmNormal}
\newcommand{\MQMobileLaneGap}{\MQRhythmLoose}

% Shared structural relationships used instead of figure-local arithmetic.
\pgfmathsetlengthmacro{\MQSectionSeparation}{8/3*\MQLayoutRhythm}
\pgfmathsetlengthmacro{\MQSectionToControlGap}{3/2*\MQLayoutRhythm}

% Canonical connector geometry. The half-span is calculated, never stored as an
% independent physical constant.
\newcommand{\MQConnectorSpan}{\MQRhythmNormal}
\pgfmathsetlengthmacro{\MQConnectorHalfSpan}{1/2*\MQLayoutRhythm}

% Base node dimensions are layout geometry and therefore derive from R.
\pgfmathsetlengthmacro{\MQProcessMinHeight}{6/5*\MQLayoutRhythm}
\pgfmathsetlengthmacro{\MQProcessMinWidth}{3*\MQLayoutRhythm}
\pgfmathsetlengthmacro{\MQDecisionMinWidth}{13/4*\MQLayoutRhythm}
\pgfmathsetlengthmacro{\MQDecisionMinHeight}{5/3*\MQLayoutRhythm}
\newcommand{\MQNoteMinHeight}{\MQRhythmNormal}
\pgfmathsetlengthmacro{\MQDatabaseMinWidth}{3*\MQLayoutRhythm}
\pgfmathsetlengthmacro{\MQDatabaseMinHeight}{3/2*\MQLayoutRhythm}

% Semantic reusable size classes. These are centrally derived relationships,
% not aliases for historical millimetre literals in individual figures.
\pgfmathsetlengthmacro{\MQNodeWidthCompact}{4*\MQLayoutRhythm}
\pgfmathsetlengthmacro{\MQNodeWidthWide}{6*\MQLayoutRhythm}
\pgfmathsetlengthmacro{\MQNodeWidthExtraWide}{15/2*\MQLayoutRhythm}
\pgfmathsetlengthmacro{\MQNodeWidthHero}{13*\MQLayoutRhythm}
\pgfmathsetlengthmacro{\MQTextWidthWide}{23/2*\MQLayoutRhythm}

% Every figure in one responsive class uses the same publication width so that
% GitHub applies the same physical scale to typography, boxes, lines and arrows.
% These fixed canvas widths are independent hard output constraints, not R-derived.
\newcommand{\MQDesktopMaxWidth}{160mm}
\newcommand{\MQMobileMaxWidth}{100mm}

% standalone is canonical at border=8pt. The internal canvas is therefore the
% publication width minus 16pt. The helper only enlarges a narrower bounding box;
% an oversized figure remains oversized and consequently fails review visibly.
\newcommand{\MQCanonicalCanvas}[1]{%
  \path (current bounding box.center) coordinate (mqcanvascenter);
  \pgfmathsetlengthmacro{\mqcanvashalf}{0.5*(#1-16pt)}%
  \path ([xshift=-\mqcanvashalf]mqcanvascenter) -- ([xshift=\mqcanvashalf]mqcanvascenter);
}
\newcommand{\MQDesktopCanvas}{\MQCanonicalCanvas{\MQDesktopMaxWidth}}
\newcommand{\MQMobileCanvas}{\MQCanonicalCanvas{\MQMobileMaxWidth}}

% ---------------------------------------------------------------------------
% Canonical connector ports
% ---------------------------------------------------------------------------

% Low-level interpolation helpers are retained for shared-system implementation
% compatibility only. Figure sources must not supply literal attachment fractions;
% use the distributed-port primitives below when several connectors share a border.
\newcommand{\MQPortNorth}[2]{($(#1.north west)!#2!(#1.north east)$)}
\newcommand{\MQPortSouth}[2]{($(#1.south west)!#2!(#1.south east)$)}
\newcommand{\MQPortWest}[2]{($(#1.south west)!#2!(#1.north west)$)}
\newcommand{\MQPortEast}[2]{($(#1.south east)!#2!(#1.north east)$)}

% Distributed ports are derived from the current border length. For port i of n,
% the canonical relative attachment fraction is i/(n+1). Figure sources provide
% only the node, coordinate name, ordinal index and total count; they never own a
% fractional attachment value.
\newcommand{\MQDistributedPortNorth}[4]{%
  \pgfmathparse{#3/(#4+1)}%
  \edef\MQDistributedPortFraction{\pgfmathresult}%
  \path ($(#1.north west)!\MQDistributedPortFraction!(#1.north east)$) coordinate (#2);%
}
\newcommand{\MQDistributedPortSouth}[4]{%
  \pgfmathparse{#3/(#4+1)}%
  \edef\MQDistributedPortFraction{\pgfmathresult}%
  \path ($(#1.south west)!\MQDistributedPortFraction!(#1.south east)$) coordinate (#2);%
}
\newcommand{\MQDistributedPortWest}[4]{%
  \pgfmathparse{#3/(#4+1)}%
  \edef\MQDistributedPortFraction{\pgfmathresult}%
  \path ($(#1.south west)!\MQDistributedPortFraction!(#1.north west)$) coordinate (#2);%
}
\newcommand{\MQDistributedPortEast}[4]{%
  \pgfmathparse{#3/(#4+1)}%
  \edef\MQDistributedPortFraction{\pgfmathresult}%
  \path ($(#1.south east)!\MQDistributedPortFraction!(#1.north east)$) coordinate (#2);%
}

% ---------------------------------------------------------------------------
% Canonical rendering / elevation
% ---------------------------------------------------------------------------

\newcommand{\MQShadowXSubtle}{0.7pt}
\newcommand{\MQShadowYSubtle}{-0.7pt}
\newcommand{\MQShadowBlurSubtle}{1.6pt}
\newcommand{\MQShadowOpacitySubtle}{0.12}
\newcommand{\MQShadowXRaised}{1.0pt}
\newcommand{\MQShadowYRaised}{-1.0pt}
\newcommand{\MQShadowBlurRaised}{2.4pt}
\newcommand{\MQShadowOpacityRaised}{0.16}

\tikzset{
  mqtt elevated none/.style={},
  mqtt elevated subtle/.style={},
  mqtt elevated raised/.style={}
}

% ---------------------------------------------------------------------------
% Typography
% ---------------------------------------------------------------------------

\newcommand{\MQBodyFont}{\sffamily\small}
\newcommand{\MQTitleFont}{\sffamily\bfseries\large}
\newcommand{\MQSectionFont}{\sffamily\bfseries\small}
\newcommand{\MQAnnotationFont}{\sffamily\footnotesize}
\newcommand{\MQCodeFont}{\ttfamily\small}
\newcommand{\MQCodeSmallFont}{\ttfamily\footnotesize}

\newcommand{\MQDesktopTypography}{%
  \renewcommand{\MQBodyFont}{\sffamily\small}%
  \renewcommand{\MQTitleFont}{\sffamily\bfseries\large}%
  \renewcommand{\MQSectionFont}{\sffamily\bfseries\small}%
  \renewcommand{\MQAnnotationFont}{\sffamily\footnotesize}%
  \renewcommand{\MQCodeFont}{\ttfamily\small}%
  \renewcommand{\MQCodeSmallFont}{\ttfamily\footnotesize}%
}

% Mobile intentionally uses a type-size uplift; layout must be recomposed rather
% than shrinking type to make a desktop composition fit.
\newcommand{\MQMobileTypography}{%
  \renewcommand{\MQBodyFont}{\sffamily\normalsize}%
  \renewcommand{\MQTitleFont}{\sffamily\bfseries\large}%
  \renewcommand{\MQSectionFont}{\sffamily\bfseries\normalsize}%
  \renewcommand{\MQAnnotationFont}{\sffamily\small}%
  \renewcommand{\MQCodeFont}{\ttfamily\normalsize}%
  \renewcommand{\MQCodeSmallFont}{\ttfamily\small}%
}

\tikzset{
  mqtt no hyphenation/.style={execute at begin node={\hyphenpenalty=10000\exhyphenpenalty=10000\relax}},
  mqtt text/.style={mqtt no hyphenation,font=\MQBodyFont,text=snodecInk,align=center},
  mqtt text left/.style={mqtt text,align=left},
  mqtt title/.style={mqtt no hyphenation,font=\MQTitleFont,text=snodecInk,align=left},
  mqtt section title/.style={mqtt no hyphenation,font=\MQSectionFont,text=snodecBlue,align=left,anchor=west},
  mqtt annotation/.style={mqtt no hyphenation,font=\MQAnnotationFont,text=snodecMuted,align=center},
  mqtt code/.style={mqtt no hyphenation,font=\MQCodeFont,text=snodecInk},
  mqtt code small/.style={mqtt no hyphenation,font=\MQCodeSmallFont,text=snodecInk}
}

% ---------------------------------------------------------------------------
% Figure presets
% ---------------------------------------------------------------------------

\tikzset{
  mqtt desktop/.style={
    x=\MQLayoutRhythm,y=\MQLayoutRhythm,
    execute at begin picture={\MQDesktopTypography},
    execute at end picture={\MQDesktopCanvas},
    every node/.append style={transform shape=false}
  },
  mqtt mobile/.style={
    x=\MQLayoutRhythm,y=\MQLayoutRhythm,
    execute at begin picture={\MQMobileTypography},
    execute at end picture={\MQMobileCanvas},
    every node/.append style={transform shape=false}
  }
}

% ---------------------------------------------------------------------------
% Core node families
% ---------------------------------------------------------------------------

\tikzset{
  mqtt node compact/.style={minimum width=\MQNodeWidthCompact},
  mqtt node wide/.style={minimum width=\MQNodeWidthWide},
  mqtt node extra wide/.style={minimum width=\MQNodeWidthExtraWide},
  mqtt node hero/.style={minimum width=\MQNodeWidthHero},
  mqtt text wide/.style={text width=\MQTextWidthWide},
  mqtt process/.style={mqtt text,draw=snodecRule!85,fill=white,rounded corners=\MQRadiusM,line width=\MQLineWidthPrimary,minimum height=\MQProcessMinHeight,minimum width=\MQProcessMinWidth,inner xsep=\MQNodePadX,inner ysep=\MQNodePadY},
  mqtt application/.style={mqtt text,draw=snodecGreen!78,fill=snodecGreenSoft,rounded corners=\MQRadiusM,line width=\MQLineWidthPrimary,minimum height=\MQProcessMinHeight,minimum width=\MQProcessMinWidth,inner xsep=\MQNodePadX,inner ysep=\MQNodePadY},
  mqtt external/.style={mqtt process,fill=snodecGraySoft,draw=snodecRule!85},
  mqtt state/.style={mqtt process,fill=snodecBlueSoft,draw=snodecBlue!78},
  mqtt decision/.style={mqtt text,diamond,aspect=1.75,draw=snodecOrange!80,fill=snodecOrangeSoft,line width=\MQLineWidthPrimary,inner sep=\MQDecisionInnerPad,minimum width=\MQDecisionMinWidth,minimum height=\MQDecisionMinHeight},
  mqtt success/.style={mqtt process,fill=snodecGreenSoft,draw=snodecGreen!78},
  mqtt error/.style={mqtt process,fill=snodecRedSoft,draw=snodecRed!82},
  mqtt control/.style={mqtt process,fill=snodecBlueSoft,draw=snodecBlue!78},
  mqtt warning/.style={mqtt process,fill=snodecOrangeSoft,draw=snodecOrange!80},
  mqtt note/.style={mqtt text,draw=snodecRule!85,fill=snodecGraySoft,rounded corners=\MQRadiusM,line width=\MQLineWidthSecondary,font=\MQAnnotationFont,minimum height=\MQNoteMinHeight,inner xsep=\MQNodePadX,inner ysep=\MQNodePadY},
  mqtt config/.style={mqtt process,fill=snodecGraySoft,draw=snodecRule!85},
  mqtt data/.style={mqtt process,fill=white,draw=snodecRule!85},
  mqtt database/.style={mqtt text,cylinder,shape border rotate=90,aspect=.25,draw=snodecRule!85,fill=snodecGraySoft,line width=\MQLineWidthPrimary,minimum width=\MQDatabaseMinWidth,minimum height=\MQDatabaseMinHeight,inner xsep=\MQNodePadX,inner ysep=\MQNodePadY}
}

% ---------------------------------------------------------------------------
% Semantic containers
% ---------------------------------------------------------------------------

\tikzset{
  mqtt figure frame/.style={draw=snodecRule!82,fill=white,rounded corners=\MQRadiusL,line width=\MQLineWidthSecondary,inner xsep=\MQGroupPad,inner ysep=\MQGroupPad},
  mqtt lane/.style={draw=snodecRule!82,fill=none,dashed,rounded corners=\MQRadiusL,line width=\MQLineWidthSecondary,inner xsep=\MQLanePadX,inner ysep=\MQLanePadY},
  mqtt group/.style={draw=snodecRule!82,fill=none,rounded corners=\MQRadiusL,line width=\MQLineWidthSecondary,inner sep=\MQGroupPad},
  mqtt data plane/.style={draw=snodecBlue!78,fill=snodecBlueSoft,rounded corners=\MQRadiusL,line width=\MQLineWidthSecondary,inner sep=\MQGroupPad},
  mqtt admin plane/.style={draw=snodecRule!85,fill=snodecGraySoft,rounded corners=\MQRadiusL,line width=\MQLineWidthSecondary,inner sep=\MQGroupPad},
  mqtt trust boundary/.style={draw=snodecRule!82,fill=none,dashed,rounded corners=\MQRadiusL,line width=\MQLineWidthSecondary,inner sep=\MQTrustPad},
  mqtt layout fit/.style={inner sep=0pt,outer sep=0pt}
}

% ---------------------------------------------------------------------------
% Arrow and connector semantics
% ---------------------------------------------------------------------------

\tikzset{
  mqtt connector/.style={line cap=round,line join=round},
  mqtt flow/.style={mqtt connector,-{Latex[length=\MQArrowPrimaryLength,width=\MQArrowPrimaryWidth]},draw=snodecMuted!76,line width=\MQLineWidthPrimary},
  mqtt data flow/.style={mqtt connector,-{Latex[length=\MQArrowPrimaryLength,width=\MQArrowPrimaryWidth]},draw=snodecBlue!70,line width=\MQLineWidthPrimary},
  mqtt control flow/.style={mqtt connector,-{Latex[length=\MQArrowPrimaryLength,width=\MQArrowPrimaryWidth]},draw=snodecBlue!70,line width=\MQLineWidthPrimary},
  mqtt observation/.style={mqtt connector,-{Latex[length=\MQArrowSecondaryLength,width=\MQArrowSecondaryWidth]},draw=snodecMuted!72,line width=\MQLineWidthSecondary,dashed},
  mqtt handoff/.style={mqtt connector,-{Latex[length=\MQArrowPrimaryLength,width=\MQArrowPrimaryWidth]},draw=snodecBlue!70,line width=\MQLineWidthPrimary},
  mqtt association/.style={mqtt connector,draw=snodecMuted!68,line width=\MQLineWidthSecondary},
  mqtt contains/.style={mqtt connector,draw=snodecMuted!68,line width=\MQLineWidthSecondary},
  mqtt dependency/.style={mqtt connector,-{Latex[length=\MQArrowSecondaryLength,width=\MQArrowSecondaryWidth]},draw=snodecMuted!68,line width=\MQLineWidthSecondary,densely dotted}
}

% ---------------------------------------------------------------------------
% Labels and evidence
% ---------------------------------------------------------------------------

\tikzset{
  mqtt branch label/.style={mqtt no hyphenation,font=\MQAnnotationFont\bfseries,text=snodecInk,fill=white,inner sep=\MQLabelInnerPad},
  mqtt edge label/.style={mqtt no hyphenation,font=\MQAnnotationFont,text=snodecMuted,fill=white,inner sep=\MQLabelInnerPad},
  mqtt edge label above/.style={mqtt edge label,above=\MQLabelGap},
  mqtt edge label below/.style={mqtt edge label,below=\MQLabelGap},
  mqtt edge label left/.style={mqtt edge label,left=\MQLabelGap},
  mqtt edge label right/.style={mqtt edge label,right=\MQLabelGap},
  mqtt branch label above/.style={mqtt branch label,above=\MQLabelGap},
  mqtt branch label below/.style={mqtt branch label,below=\MQLabelGap},
  mqtt branch label left/.style={mqtt branch label,left=\MQLabelGap},
  mqtt branch label right/.style={mqtt branch label,right=\MQLabelGap},
  mqtt topic label/.style={mqtt no hyphenation,font=\MQCodeSmallFont,text=snodecInk,fill=white,inner sep=\MQLabelInnerPad},
  mqtt route label/.style={mqtt no hyphenation,font=\MQCodeSmallFont,text=snodecInk,fill=white,inner sep=\MQLabelInnerPad},
  mqtt evidence runtime/.style={mqtt no hyphenation,font=\MQAnnotationFont\bfseries,text=snodecGreen!78,draw=snodecGreen!78,fill=snodecGreenSoft,rounded corners=\MQRadiusM,line width=\MQLineWidthSecondary,inner xsep=\MQEvidencePadX,inner ysep=\MQEvidencePadY},
  mqtt evidence source/.style={mqtt no hyphenation,font=\MQAnnotationFont\bfseries,text=snodecBlue!78,draw=snodecBlue!78,fill=snodecBlueSoft,rounded corners=\MQRadiusM,line width=\MQLineWidthSecondary,inner xsep=\MQEvidencePadX,inner ysep=\MQEvidencePadY}
}

% ---------------------------------------------------------------------------
% Canonical publication rules
% ---------------------------------------------------------------------------
%
% 1. Figure .tex + this style file are canonical.
% 2. Generated SVGs are never hand-edited or committed as canonical source.
% 3. SVG is the publication format for technical diagrams.
% 4. PNG is reserved for genuine runtime/UI evidence or raster-only input.
% 5. Desktop/mobile may share semantics but use independently art-directed geometry.
% 6. No arrow without real directional semantics.
% 7. Dashed arrows mean observation/secondary relationship, not primary flow.
% 8. Containers denote real process/runtime/trust/configuration/lifecycle boundaries.
% 9. Color never carries the only semantic distinction.
% 10. All spatial layout tokens derive from the one canonical rhythm R; figure sources never own physical layout values or local multipliers.
% 11. Horizontal sequences use the canonical horizontal rhythm; avoid arbitrary x-spacing.
% 12. Rounded corners follow the semantic radius scale; avoid one-off radii.
% 13. Flat rendering is authoritative; do not introduce decorative elevation/shadows.
% 14. Prefer explicit anchors/ports and orthogonal routing wherever boxes are connected.
% 15. Preserve readable text size at actual GitHub desktop/mobile rendering widths.
% 16. Every directional connector is a single path with its arrowhead attached to that path; detached or separately placed arrowheads are forbidden.
% 17. Connectors leave and enter box borders orthogonally at 90 degrees. Use the middle of the relevant border by default.
% 18. If several connectors share one border, use the canonical distributed-port primitives. Port i of n is derived from the current border geometry at i/(n+1); figure sources never own attachment fractions or distances.
% 19. Do not use unmotivated diagonal connectors. Use orthogonal Manhattan routing between ports unless a diagonal has a specific semantic or geometric justification.
% 20. Use mqtt figure frame for the neutral whole-figure surface; draw it before semantic lanes/groups/planes. Never let an enclosing fill overpaint inner semantics.
% 21. Desktop and mobile typography are selected by their presets; do not override shared font sizes ad hoc to compensate for an oversized composition.
% 22. \MQDesktopMaxWidth and \MQMobileMaxWidth are fixed publication canvas widths and hard composition limits. Recompose a figure that exceeds them instead of scaling it down.
% 23. Automatic word hyphenation is forbidden in diagram labels, notes, routes, and code tokens.
% 24. mqtt success is reserved for a real successful outcome. Do not use green merely for a destination, final stage, action, or emphasis.
% 25. Ownership/containment uses mqtt contains or mqtt association, never a flow arrow unless real directional flow also exists.
% 26. Repeated label placement uses the canonical \MQLabelGap through shared directional label styles; labels do not sit directly on connector shafts.
